% !TEX program = xelatex
\documentclass[UTF8,a4paper,11pt]{ctexart}

\usepackage[a4paper,margin=2.2cm]{geometry}
\usepackage{booktabs}
\usepackage{longtable}
\usepackage{array}
\usepackage{xcolor}
\usepackage{listings}
\usepackage{enumitem}
\usepackage{fancyhdr}
\usepackage{xurl}
\usepackage{hyperref}

\setlength{\headheight}{14pt}
\renewcommand{\UrlFont}{\ttfamily\small}

\hypersetup{
  colorlinks=true,
  linkcolor=black,
  urlcolor=blue!50!black,
  pdftitle={ChargeHub Socket 网络编程详解},
  pdfauthor={ChargeHub}
}

\pagestyle{fancy}
\fancyhf{}
\fancyhead[L]{ChargeHub}
\fancyhead[R]{Socket 网络编程详解}
\fancyfoot[C]{\thepage}
\renewcommand{\headrulewidth}{0.4pt}

\lstset{
  basicstyle=\ttfamily\small,
  breaklines=true,
  columns=fullflexible,
  frame=single,
  framesep=4pt,
  xleftmargin=6pt,
  xrightmargin=6pt,
  backgroundcolor=\color{gray!6},
  showstringspaces=false,
  language=C++
}

\title{ChargeHub Socket 网络编程详解\\[0.35em]
\large 从网络基础到本项目的 TCP、HTTP 与外部接口}
\author{对照作业目录当前源码}
\date{2026 年 9 月}

\begin{document}
\maketitle

\begin{abstract}
\noindent
本文按教科书顺序讲：先 IP/端口/TCP，再 Socket 与 Qt 封装，再讲本项目自己定的长度前缀 JSON，
最后讲 HTTP/HTTPS 大屏和腾讯/高德等出站接口。读完应能回答：为什么用户端不直连数据库、
粘包怎么切、\texttt{readyRead} 里为什么必须 \texttt{readAll}、天气请求为什么不走 8888。
不粘贴 Key / SK / 密码。字段表见 \texttt{protocol/messages.md}。
\end{abstract}

\tableofcontents
\newpage

\section{先建立地图：本项目里有几条「网」}

ChargeHub 不是单机计算器，是\textbf{至少两个进程}通过网络说话：

\begin{center}
\begin{tabular}{lll}
\toprule
通道 & 传输层 & 本项目用法 \\
\midrule
用户端 $\leftrightarrow$ 管理端 & TCP 端口 8888 & 自己设计的长度前缀 JSON（Socket 编程主体） \\
浏览器 $\leftrightarrow$ 大屏 & TCP 上的 HTTP 5000 & Flask 只读 JSON，标准 Web \\
用户端/管理端 $\to$ 互联网 & TCP 上的 HTTP/HTTPS & 腾讯、高德瓦片、Open-Meteo 等 \\
\bottomrule
\end{tabular}
\end{center}

三条通道都建立在 \textbf{TCP} 上，但\textbf{应用层协议不同}：
8888 是我们自己的二进制帧；5000 和腾讯是 HTTP 文本/JSON。
操作系统只保证把字节按序送到，\textbf{不理解} JSON 或 ``LOGIN''。
把字节解释成业务，是应用层（我们的 \texttt{Protocol}、Flask、浏览器）的事。

\section{网络基础}

\subsection{IP 地址：找到那台电脑}

IP 地址标识一台主机在某个网络里的位置。本项目常见三种：

\begin{description}[style=nextline,leftmargin=1.2em]
  \item[\texttt{127.0.0.1}]
    本机回环。用户端和管理端在\textbf{同一台电脑}时填它。数据包不出网卡。
  \item[局域网地址，如 \texttt{192.168.1.8}、WSL 的 \texttt{172.29.x.x}]
    同一 Wi-Fi / 同一虚拟网里别人能 ping 到的地址。管理端底栏显示的就是这种。
    组员用户端必须填\textbf{值班那台}的局域网 IP，填自己的 127.0.0.1 会连到自己机器（往往没开管理端）。
  \item[公网 IP]
    校园网/家用路由后面通常没有真正的公网监听。本作业按局域网联调，不要求公网映射。
\end{description}

IPv4 是 32 位，点分十进制四段。本作业只使用 IPv4（\texttt{QHostAddress::Any} 实际听 IPv4 的 0.0.0.0）。

\subsection{端口：找到那台电脑上的哪个程序}

一台电脑可以同时跑浏览器、管理端、大屏。端口把它们分开：

\begin{itemize}[leftmargin=1.4em]
  \item 端口是 16 位整数，范围 0--65535。
  \item \textbf{服务端}绑定一个约定端口等别人来：本项目管理端 \textbf{8888}，大屏 \textbf{5000}。
  \item \textbf{客户端}通常由操作系统分配临时端口（例如 49152 以上），你在登录页填的是\textbf{对方}的端口，不是自己的。
\end{itemize}

所以完整地址是 \texttt{IP:端口}。登录页 \texttt{127.0.0.1:8888} 的意思是：
到本机上那个正在 \texttt{listen(8888)} 的进程。浏览器打开 \texttt{http://127.0.0.1:5000} 则是另一个进程。
\textbf{不要用浏览器打开 8888}：浏览器讲 HTTP，管理端 8888 讲我们自己的帧，对不上。

\subsection{客户端与服务器}

\begin{itemize}[leftmargin=1.4em]
  \item \textbf{服务器}先启动、先 \texttt{listen}，被动等待。管理端必须先开。
  \item \textbf{客户端}后启动、主动 \texttt{connect}。用户端、浏览器、烟测脚本都是客户端。
\end{itemize}

同一台机器可以既是 8888 的服务器，又是腾讯 HTTPS 的客户端：管理端相对用户端是服务器，
相对 \texttt{apis.map.qq.com} 是客户端。角色取决于谁先听、谁去连。

\subsection{分层模型（够用即可）}

不必背全 OSI 七层。记住四层就够把本项目对上：

\begin{center}
\begin{tabular}{lll}
\toprule
层 & 本项目对应 & 你改代码时碰不碰 \\
\midrule
应用层 & JSON / HTTP / 腾讯 WebService & 碰：\texttt{type}、URL、Flask 路由 \\
传输层 & TCP（可靠字节流） & 间接碰：端口、\texttt{QTcpSocket} \\
网络层 & IP 路由、127.0.0.1 vs 局域网 & 联调时填对地址 \\
链路/物理 & 网卡、Wi-Fi、虚拟网桥 & 一般不改代码 \\
\bottomrule
\end{tabular}
\end{center}

Qt 的 \texttt{QTcpSocket} 帮你走完传输层；\texttt{Protocol} 和 Flask 才是应用层。

\section{TCP 为什么适合充电业务}

\subsection{TCP 提供什么}

TCP（Transmission Control Protocol）在两台主机之间提供：

\begin{itemize}[leftmargin=1.4em]
  \item \textbf{面向连接}：先三次握手，双方确认后再传数据。对应 Qt 的 \texttt{connected} 信号。
  \item \textbf{可靠}：丢包会重传，重复会丢弃，乱序会排好。应用层看到的是按发送顺序的字节。
  \item \textbf{双向字节流}：两边都能随时 \texttt{write}。所以服务端能主动推 \texttt{PUSH\_CHARGE}。
\end{itemize}

三次握手只需记住：客户端 \texttt{connect} 成功，说明连接已建立，可以发 LOGIN。
四次挥手对应断开：用户端关掉窗口，管理端收到 \texttt{disconnected}。

\subsection{为什么不用 UDP}

UDP 不握手、不保证到达和顺序，适合游戏位置广播、视频。充电要：
开充必须到达、结算不能乱序、推送最好别丢到「界面永远停住」（我们另外用轮询兜底）。
所以业务通道选 TCP。地图瓦片虽然也是 TCP 上的 HTTP，失败只换下一张 URL，不改账。

\subsection{TCP 不提供「消息」}

这是 Socket 编程里最容易错的一点：

\begin{quote}
TCP 给你的是\textbf{字节流}，不是「一封信」。一次 \texttt{write} 可能被拆成两次 \texttt{readyRead}；
两次 \texttt{write} 也可能粘在一次 \texttt{read} 里。这叫\textbf{半包}和\textbf{粘包}。
\end{quote}

如果直接 \texttt{readAll()} 当一条 JSON，偶发会解析失败或把两单粘成一次。
所以必须自己在字节流上划边界。本项目的划法：前面 4 字节告诉你后面 JSON 有多长。

\section{字节序、UTF-8 与长度前缀}

\subsection{大端与小端}

多字节整数在内存里有两种排法。网络约定用\textbf{大端}（高位字节先发出），又称网络字节序。
x86/Windows/WSL 本机一般是小端，所以写长度时要转换：

\begin{lstlisting}
qToBigEndian<quint32>(body.size(), out.data());   // 发送
qFromBigEndian<quint32>(buf.constData());           // 接收
\end{lstlisting}

Python 烟测与此完全一致：\texttt{struct.pack(">I", len(body))}，\texttt{>} 就是大端。

\subsection{UTF-8 JSON}

长度统计的是\textbf{字节}，不是汉字个数。「北京」UTF-8 是 6 字节。
\texttt{QJsonDocument::Compact} 去掉空格，两端必须同一套编码，否则长度对、内容却解不出来。

\subsection{一帧的样子}

\begin{lstlisting}[language={}]
+------------+---------------------------+
| 4 bytes    | N bytes UTF-8 JSON        |
| big-endian | {"type":"LOGIN", ...}     |
| length = N |                           |
+------------+---------------------------+
\end{lstlisting}

非法长度（0 或大于 1\,MB）时，接收端清空缓冲：宁可变一次丢包，也不要永远等一个假长度。

\section{Socket 是什么：从 BSD 到 Qt}

\subsection{套接字 = 编程入口}

Socket 是操作系统提供的「网线插头」。服务端、客户端各拿一个，连上之后就能读写字节。
经典 BSD 接口与本项目 Qt 的对应：

\begin{center}
\begin{tabular}{lll}
\toprule
BSD / POSIX & Qt（本项目） & 发生在哪 \\
\midrule
\texttt{socket()} & 构造 \texttt{QTcpSocket} / \texttt{QTcpServer} & 客户端 / 服务端 \\
\texttt{bind()+listen()} & \texttt{QTcpServer::listen(Any, 8888)} & \texttt{adminserver/src/main.cpp} \\
\texttt{accept()} & \texttt{incomingConnection} 里 \texttt{setSocketDescriptor} & \texttt{tcpserver.cpp} \\
\texttt{connect()} & \texttt{QTcpSocket::connectToHost} & \texttt{client.cpp} \\
\texttt{send()/write()} & \texttt{socket->write(Protocol::pack(...))} & 两端 \\
\texttt{recv()/read()} & \texttt{readyRead} 里 \texttt{readAll()} & 两端 \\
\texttt{close()} & 窗口关 / \texttt{abort()} / \texttt{disconnected} & 两端 \\
\bottomrule
\end{tabular}
\end{center}

一条 TCP 连接可以用五元组标识：协议、本机 IP、本机端口、对端 IP、对端端口。
管理端一个 \texttt{listen} 套接字可以 \texttt{accept} 出很多根已连接套接字，每个用户端一根。

\subsection{阻塞 vs 事件驱动}

教科书里 \texttt{recv()} 会卡住线程直到有数据。Qt GUI 不能卡死界面，所以用\textbf{事件循环}：

\begin{itemize}[leftmargin=1.4em]
  \item 把套接字交给 Qt，有数据时发 \texttt{readyRead} 信号。
  \item 槽函数里一次性 \texttt{readAll()}，喂给 \texttt{Protocol::append}。
  \item 再 \texttt{while (hasPacket())} 取出完整 JSON。半包留在 \texttt{Protocol} 的缓冲里，等下次信号。
\end{itemize}

这就是「网络编程」在桌面程序里的标准形态：\textbf{非阻塞 + 回调/信号}。

\section{本项目的 TCP 编程（8888）}

\subsection{服务端：先听再接}

\texttt{adminserver/src/main.cpp} 在管理员登录框弹出\textbf{之前}就：

\begin{lstlisting}
server.listen(QHostAddress::Any, 8888);
\end{lstlisting}

\texttt{Any} 表示所有网卡（0.0.0.0），所以局域网里别人才能连进来。
若只听 127.0.0.1，组员跨机器会失败。听失败通常是第二份管理端还在抢端口。

每个新连接进入 \texttt{TcpServer::incomingConnection}：

\begin{enumerate}[leftmargin=1.6em]
  \item 用内核给的描述符构造一根 \texttt{QTcpSocket}。
  \item \textbf{每根线一个} \texttt{Protocol} 实例（缓冲不能混用）。
  \item \texttt{readyRead}：拆包 $\to$ \texttt{Dispatch::handle} $\to$ \texttt{pack} 写回。
        处理是同步的：这一帧的 SQL 跑完才回包。作业并发量小，不用线程池。
  \item \texttt{disconnected}：摘掉绑定；该用户已无其它线则 60 秒后释放桩。
\end{enumerate}

\subsection{客户端：拨号再发业务}

\texttt{Client::connectTo} 先 \texttt{abort} 旧连接再 \texttt{connectToHost(host, port)}，保证一个 \texttt{Client} 最多一根 TCP。
只有 \texttt{ConnectedState} 才 \texttt{write}，否则只报「未连接到服务器」，避免写半包。

\texttt{request} 组信封后 \texttt{Protocol::pack} 再 \texttt{write+flush}。
\texttt{seq} 在本连接上自增。服务端推送的 \texttt{PUSH\_CHARGE} 的 seq 固定为 0，与请求对不上号，界面按 \texttt{type} 区分。

\subsection{接收循环为什么必须写成这样}

\begin{lstlisting}
codec_.append(sock_.readAll());
while (codec_.hasPacket())
    emit responded(codec_.nextPacket());
\end{lstlisting}

\begin{itemize}[leftmargin=1.4em]
  \item 必须 \texttt{readAll}：内核缓冲里可能已经粘了多帧。
  \item 必须 \texttt{while}：一次信号可能拆出 0、1、多张 JSON。
  \item 半包不报错：等下一轮 \texttt{readyRead}。
\end{itemize}

\section{应用层协议：我们自己的「HTTP」}

TCP 只运字节。8888 上的语法是我们定的，相当于一个极简 HTTP。

\subsection{请求信封}

\begin{lstlisting}[language={}]
{ "type": "LOGIN", "seq": 1, "role": "user", "token": "", "data": { ... } }
\end{lstlisting}

\begin{itemize}[leftmargin=1.4em]
  \item \texttt{type}：方法名（类似 HTTP 的路径 + 动词合在一起）。
  \item \texttt{seq}：这一帧的序号，方便日志；写操作还拿它做幂等键的一部分。
  \item \texttt{token}：登录后的身份，类似 Cookie / Authorization。
  \item \texttt{data}：方法参数。
\end{itemize}

\subsection{响应信封}

\begin{lstlisting}[language={}]
{ "type": "LOGIN", "seq": 1, "code": 0, "message": "登录成功", "data": { ... } }
\end{lstlisting}

\texttt{code==0} 成功。失败时 \texttt{data} 为空，避免界面把半成品当成功。
常见码与 HTTP 故意对齐：400 参数、401 未登录、403 禁止、404 没有、409 冲突。

\subsection{请求--响应，以及服务器推送}

多数 type 是一问一答：客户端发，服务端立刻回同一 \texttt{seq}。
充电中额外有\textbf{服务器推送}：\texttt{TcpServer} 每 5 秒对在线用户 \texttt{write} 一帧 \texttt{PUSH\_CHARGE}，客户端没有先发这个 type。
这是 TCP 双向流才能做的事；若改成短连接 HTTP，就要改成客户端轮询或 SSE/WebSocket。

用户端仍每 1 秒自己问 \texttt{CHARGE\_STATUS}：推送丢了或旧客户端忽略推送时，界面还能动。

\subsection{鉴权发生在哪一层}

\texttt{Protocol} 不看 token。\texttt{TcpServer} 只搬运。
\texttt{RequestDispatcher} 查路由表：\texttt{LOGIN}/\texttt{REGISTER} 公开；其余先 \texttt{SessionService::authenticate}。
Token 是 UUID 字符串，落 \texttt{session} 表，30 分钟滑动过期。这是应用层会话，不是 TLS 证书。

\subsection{幂等：网络编程里的业务问题}

用户连点「开始充电」，可能发出两帧相同 \texttt{type} 和 \texttt{seq}。
TCP 保证都到达，业务层必须认「这是同一操作」。
键 \texttt{token|type|seq}，8 秒内返回上次成功响应。这不是 TCP 重传，是应用层防重复提交。

\section{HTTP 与 HTTPS：另一套应用层}

\subsection{HTTP 报文长什么样}

浏览器和大屏、Qt 访问腾讯，用的是 HTTP：一行方法 + 路径，然后头部，空行，可选正文。例如：

\begin{lstlisting}[language={}]
GET /api/overview HTTP/1.1
Host: 127.0.0.1:5000
User-Agent: ChargeHub/1.0 (campus training)

\end{lstlisting}

响应第一行是状态码 \texttt{200 OK}，正文常是 JSON。
Flask 的 \texttt{@app.get("/api/overview")} 就是在应用层注册这条路径。

HTTPS = HTTP 跑在 TLS 之上（传输仍是 TCP，多了加密握手）。
腾讯 \texttt{https://apis.map.qq.com} 由 Qt 网络库做证书校验；业务代码只看到 URL 和 JSON。
本作业 8888 \textbf{没有} TLS，局域网作业可接受；不要把 token 打到公共 Wi-Fi 的演示里当安全课结论。

\subsection{大屏：HTTP 只读}

\texttt{dashboard/app.py} 监听 \texttt{0.0.0.0:5000}。三个 API 全是 GET，SQL 只有 SELECT。
没有 token：能访问 5000 的人只能看。这是有意的产品边界，不是「忘了做登录」。

管理端用本机 TCP 探 5000 是否在听，没有则拉起 Flask。这是「用 Socket 探端口」的典型用法，不是用户协议。

\subsection{出站 WebService}

\texttt{TencentApi::signedUrl} 拼
\texttt{https://apis.map.qq.com} + 路径 + 排序查询串；有 SK 时对
\texttt{path + "?" + query + SK} 做 MD5，追加 \texttt{sig}。这是腾讯文档规定的签名，不是 Socket 帧。

用户端 \texttt{QNetworkAccessManager::get} 是\textbf{异步}：发完立刻回界面，完成时槽函数读 \texttt{QNetworkReply}。
管理端预测用 \texttt{fetchUrl}：临时 \texttt{QEventLoop} 等到超时或完成，会卡住调用线程约数秒，所以不要放在拖地图的热路径里。

瓦片 URL 是普通 HTTP GET 图片。失败换列表里下一张（高德 $\to$ 腾讯）。Referer 头是防盗链要求，属于 HTTP 头，不是 TCP 选项。

\subsection{DNS 与超时}

\texttt{connectToHost("apis.map.qq.com")} 或 \texttt{QUrl} 会先 DNS 再连 443。
校园网 DNS 慢或墙掉 OSM 时，本项目已改走国内瓦片和腾讯，避免依赖 \texttt{router.project-osrm.org}。
所有出站带约 8 秒传输超时，避免界面永久转圈。

\section{连接生命周期与故障}

\subsection{正常关闭}

用户关掉用户端 $\to$ TCP 断开 $\to$ 服务端 \texttt{disconnected}。
若该用户没有其它用户端还连着，开始 60 秒计时；到点仍离线，充电中订单改待结算、桩回闲置。
这是应用层超时，不是内核 TCP keepalive 的默认 2 小时。

\subsection{异常断电}

对端崩溃可能既不挥手也不立刻被内核发现。所以有 \texttt{HEARTBEAT}（应用层保活并滑动 token）
和 60 秒释放。仅靠 TCP 保活对「人走了桩还占用」反应太慢。

\subsection{重连}

\texttt{connectTo} 会 \texttt{abort} 旧线。Token 若仍在 30 分钟内，不必重新 LOGIN（「记住我」是本机设置自动再发 LOGIN，两套机制）。
幂等缓存不跨进程持久化：管理端重启后连点保护从头算，token 仍可从库恢复。

\subsection{一人两根线}

同一账号两台电脑登录：两根 TCP，两份推送。释放桩看「是不是最后一根也断了」。
这是会话设计，不是 TCP 本身的功能。

\section{地址怎么填：联调时的网络}

\begin{center}
\begin{tabular}{lp{8.8cm}}
\toprule
场景 & 用户端填写 \\
\midrule
本机自测 & \texttt{127.0.0.1:8888} \\
组员连值班电脑 & 值班管理端底栏的 \texttt{IP:8888} \\
WSL 里听、Windows 里测 & 有时要用发行版 IP，不能想当然 127.0.0.1 一定转发 \\
防火墙 & 服务器需放行 8888、5000；客户端做出站连接即可 \\
\bottomrule
\end{tabular}
\end{center}

虚拟机网卡要能互通（桥接或同一虚拟网）。NAT 只出站时，别人连不进你的 8888。
这是网络层/部署问题，改 JSON 字段解决不了。

\section{对照源码：建议阅读顺序}

\begin{enumerate}[leftmargin=1.6em]
  \item \texttt{common/protocol.cpp}：长度、大端、粘包。可与 \texttt{scripts/smokelocal.py} 对照。
  \item \texttt{userclient/src/client.cpp}：\texttt{connectToHost}、\texttt{write}、\texttt{readyRead}。
  \item \texttt{adminserver/src/main.cpp} 的 \texttt{listen}，再 \texttt{tcpserver.cpp} 的 \texttt{incomingConnection}。
  \item \texttt{transport/requestdispatcher.cpp}：鉴权与幂等（应用层）。
  \item \texttt{services/sessionservice.cpp}：token 与 30 分钟。
  \item \texttt{dashboard/app.py}：HTTP GET 只读。
  \item \texttt{common/tencentapi.cpp}：HTTPS URL、签名、同步/异步 GET、失败链路。
\end{enumerate}

动手：先开管理端，再运行

\begin{lstlisting}[language=Python]
python scripts/smokelocal.py
\end{lstlisting}

它用标准库 \texttt{socket.create\_connection} 走完全相同的帧。能跑通，说明你理解的不是 Qt 特有魔法，而是 TCP + 长度前缀 + JSON。

\section{把知识收成一张表}

\begin{center}
\begin{tabular}{lll}
\toprule
概念 & 本项目落点 & 常见误解 \\
\midrule
IP:端口 & \texttt{127.0.0.1:8888} / 底栏 IP & 用浏览器打开 8888 \\
TCP 字节流 & \texttt{Protocol} 切帧 & 一次 read 等于一条消息 \\
大端长度 & \texttt{qToBigEndian} / \texttt{>I} & 按本机小端写长度 \\
accept 多连接 & 每用户一个 \texttt{QTcpSocket} & 一个 listen 只能服务一个人 \\
事件驱动 & \texttt{readyRead} & 在 GUI 线程里阻塞 \texttt{recv} \\
应用层方法 & JSON \texttt{type} & 以为 TCP 知道 LOGIN \\
会话 & token + session 表 & 把 TCP 连接当成永久身份 \\
推送 & 服务端主动 \texttt{write} & 以为只有客户端能说话 \\
HTTP & 大屏 5000、腾讯 HTTPS & 把 8888 当成网站 \\
失败隔离 & 天气失败不改订单 & 外网挂了就不能结算 \\
\bottomrule
\end{tabular}
\end{center}

\vspace{1em}
\noindent\textit{答辩怎么讲见《最终答辩-网络通信与外部接口》；逐文件对照见《网络通信与外部接口说明》。}

\end{document}
